% !PRINTED 581
\weberpage{581}
%<<CONT
or finally, by (2) and (3),
\weq{12}{-27\gamma = nA + 3n - 1.}

From this relation it follows, since $\gamma$ is an integer,
\weq{13}{A \equiv 1\ (\mathrm{mod}\ 3)}
and from this $A^3 \equiv 1\ (\mathrm{mod}\ 9)$. Now by (2), (3) and (6)
\weq{14}{n = a^2 - ab + b^2}
or
\weq{15}{4n = A^2 + 3b^2}
and if we substitute this value in (12),
\[
-4\cdot 27\gamma = A^3 + 3Ab^2 + 3A^2 + 9b^2 - 4,
\]
and from this
\[
3b^2 \equiv 0\ (\mathrm{mod}\ 9).
\]

Thus $b$ is divisible by 3, and if we set $b = 3B$,
it follows from (15) that
\weq{16}{4n = A^2 + 27B^2.}

We can also easily form the square root of the dis\-criminant, namely
\weq{17}{\sqrt{D} = (\eta - \eta_1)(\eta - \eta_2)(\eta_1 - \eta_2) = s - s',}
if we subtract (10) and (11) from one another. From this one
obtains
\weq{18}{\sqrt{D} = nB.}

The cubic equation whose roots are the three quantities $\eta$
therefore has, as indeed already follows from what was remarked at the
beginning of \S.\,171, three real roots. It is simplified
if one sets
\[
3\eta + 1 = \xi
\]
and then yields
\weq{19}{\xi^3 - 3n\xi - nA = 0.}

If one substitutes on the left-hand side for $\xi$ in succession
\weq{20}{-2\sqrt{n}, \quad -\sqrt{n}, \quad +\sqrt{n}, \quad +2\sqrt{n},}
there result the values
\[
-n(2\sqrt{n} + A), \quad +n(2\sqrt{n} - A), \quad -n(2\sqrt{n} + A),
\]
\[
+n(2\sqrt{n} - A),
\]
which, because by (16) $A$ is absolutely smaller than $2\sqrt{n}$, have alternating
signs. Hence in each of the three intervals bounded by the
values (20) there lies \gesperrt{one} root of equation (19).
